\subsection{Самосопряжённые операторы}

\begin{definition}
Пусть $H(\C)$~--- гильбертово пространство над полем $\C$,
$A \in \L(H) : A^* = A$, то есть $(Ax, y) = (x, Ay)\ \forall x, y \in H$.
Такой оператор $A$ называется \textit{самосопряжённым}.
\end{definition}

\begin{exercise}[теорема Хеллингера-Тёплица]
 $A: H \to H$ - линейный и симметричный, тогда $A$ - ограничен. (подсказки к решению: Рисс-Фреше, Банах-Штейнгауз)
\end{exercise}

\begin{definition}
    Если $A^*A=AA^*$, то оператор называется \textit{нормальным}
\end{definition}



\textbf{Основная идея:} $A: H \to H$, будем вместо него рассматривать квадратичную форму:
\[ K: H \to \C \]
\[ K(x)=(Ax, x) \]

\begin{exercise}
    Если $A \in \L(H)$ и $K(x)$ вещественно при любом $x$, то $A$~---
    самосопряжённый оператор.
\end{exercise}

% аналогия с тем, что таблицу умножения можно восстановить только по значениям на диагонали (только квадраты)
% в этом месте мы клонируем Коновалова

\begin{exercise}
$H(\C)$~--- гильбертово пространства, $A \in \L(H)$,
$(Ax, x) = 0\ \forall x$. Тогда $A = 0$.
Для доказательства нужно использовать 3 следствие из теоремы Хана-Банаха.

Вторая часть: для $H(\R)$ это неверно.
\end{exercise}

% Нас различают паспорта
% начинаем новую жизнь

\begin{theorem}[11.1]\label{11.1}
$H(\C)$, $A$~--- самосопряжённый оператор.
\begin{enumerate}
    \item $K(x) = (Ax, x) \in \R \ \forall x$
    \item $\lambda$~--- собственное значение оператора $A$, 
    то $\lambda \in \R$
    \item $\lambda_1 \neq \lambda_2$~--- собственные значения, 
    $e_1, e_2$~--- собственные векторы, соответствующие им. 
    Тогда $(e_1, e_2) = 0$.
\end{enumerate}
\end{theorem}

\begin{proof}
    \begin{enumerate}
        \item В силу самосопряженности $A$: $(x, Ax)=(Ax, x)=\overline{(x, Ax)}$. Число совпадает с сопряжением, значит оно вещественное.
        
        \item $\lambda$~--- собственное значение, $e$~-- собственный вектор.
        
        С другой стороны, 
        $(Ae, e) = (\lambda e, e) = \lambda (e, e)$. Но из пункта 1
        $(Ae, e) \in \R$
    
        \item $e_1, e_2, \lambda_1, \lambda_2$. 
        \[ (\lambda_1 e_1, e_2) = (Ae_1, e_2)=(e_1, Ae_2) = (e_1, \lambda_2 e_2)\]
        В силу пункта 2 $\lambda_1, \lambda_2$~--- вещественные, значит можно вынести из скалярного произведения
        \[ (\lambda_1 - \lambda_2) (e_1, e_2) = 0\]
        По условию собственные значения не равны, значит $(e_1, e_2)=0$
    \end{enumerate}
\end{proof}

\begin{theorem}[11.2]\label{11.2}
    $H(\C)$, $A$~--- самосопряжённый оператор.
    \[
    \overline{\im A_{\lambda}} \oplus_{\perp} \ker A_{\lambda} = H
    \]
\end{theorem}

\begin{proof}
    По теореме \nameref{9.2} $\overline{\im A_\lambda} \oplus \ker (A_\lambda)^*=H$
    $(A - \lambda I)^* = A^* - \overline{\lambda} I$
    % проблема: вылезет предательская лямбда с чертой
    Рассмотрим два варианта
    \begin{enumerate}
        \item $\lambda \in \R$
        
        $\lambda = \overline{\lambda}$, тогда формула автоматически верна
        \item $\lambda \not\in \R$, $\overline{\lambda} \neq \lambda$ 
        Отсюда по теореме \nameref{11.1} (пункт 2) $\lambda$ и $\overline{\lambda}$
        не являются собственным значением, а значит ядро тривиально.
        \end{enumerate}
\end{proof}

\begin{theorem}[11.3 (Критерий принадлежности числа спектру)]\label{11.3}

    $H(\C)$, $A$~--- самосопряжённый оператор.
    
    \begin{enumerate}
        \item $\lambda \in \rho(A) \lra A$~--- ограничен снизу, т.е. $\exists m > 0: \; \| A_\lambda x \| \geqslant m \| x \| \; \forall x \in H$
        
        \item $\lambda \in \sigma(A) \lra \exists \text{последовательность} \{x_n\}$, такая что 
        $\|x_n\| = 1 \ \forall n$ и $A_{\lambda} x_n \to 0$.
    \end{enumerate}
\end{theorem}
% достаточно доказать один пункт и аккуратно переписать
\begin{proof}
    Знаем, что ограниченность снизу гарантирует обращение оператора на образе
    \begin{itemize}
        \item[$\Rightarrow$:] очевидно по теореме \nameref{8.1}
        \item[$\Leftarrow$:] 
        Имеем из теоремы \nameref{11.2}: $\overline{\im A_{\lambda}} \oplus \ker A_{\lambda} = H$.
        Разумеется, ядро тривиально (внимательно посмотрим на условие ограниченности снизу). Но, как мы знаем из доказательства теоремы \nameref{8.4}, ограниченность снизу влечёт замкнутость образа. Распишем это подробнее.

        Почему из условия ограниченности снизу следует замкнутость образа?

        $\{y_n\} \in \im A_{\lambda}, y_n \to y$. Хотим доказать, что $y$ лежит в образе
        
        Раз она сходится, то она фундаментальна, то есть $\| y_{n+p}-y_n \| \to 0$. 

        \[
        \|y_{n + p} - y_n\| = \|A_{\lambda} x_{n + p}- A x_n\| = 
        \|A (x_{n + p} - x_n)\|
        \]

        Получаем, что последовательность $x_n$ фундаментальна, а значит сходится. По непрерывности оператора $y_n = A_\lambda x_n \to Ax = y$.
        Это и означает, что $y$ лежит в образе оператора $A$.
        
        Тогда получаем, что $\im A_{\lambda} = H$. Значит получили ограниченность снизу для всего $H$.
    \end{itemize}

    Переделка: $\forall m: \; m = \frac 1 n \; \exists x_n \; \| A_\lambda x_n \| < \frac{1}{n}$
\end{proof}

% в теореме 11.1 сказали что все собственные значения вещественные. как можно усилить?
\begin{theorem}[11.4]\label{11.4}
    $H(\C)$, $A$~--- самосопряжённый оператор.
    % формулировка по Андрею
    (Весь спектр) $\sigma(A) \subset \R$.
    Равносильная формулирвока: если $\lambda \not\in \R \implies \lambda \in \rho(A)$.
\end{theorem}

\begin{proof}
Идея: покажем, что если $\nu = \im \lambda \neq 0$ (мнимая часть), то $A_{\lambda}$~--- ограниченный снизу. Точнее, $\| A_\lambda x \| \geqslant |\nu| \| x \|$. $\lambda = \mu + i \nu$

% обычная арифметическая элементарщина
% TODO зачеркнуть среднюю строку в формуле
\begin{multline*}
\|A_{\lambda} x\|^2 = (A_{\lambda} x, A_{\lambda} x) = 
(A - (\mu + i \nu) x, A - (\mu + i \nu) x) = \\
= \underbrace{\|A x\|^2}_{\geqslant 0} - 
(\mu - i \nu) (A x, x) - (\mu + i \nu) (x, Ax) +
|\mu + i \nu\|^2 \|x\|^2
= (A_{\mu} x - i \nu x, A_{\mu} - i \nu x) = \\
= \|A_{\mu} x\|^2 - i \nu (x, A_{\mu} x) + i \nu (A_{\mu}x, x) +
|\nu|^2 \|x\|^2 \geqslant |\nu|^2 \|x\|^2
\end{multline*}

В последнем переходе слагаемые сократились за счет самосопряженности.
% нужно все таки понимать к чему мы стремимся: ню должен быть отделен от вещественной части. поэтому переписали абсолютно другим способом

\end{proof}

\textbf{Обозначения:} $m_+=\sup_{\|x\|=1} (Ax, x)$, $m_-=\inf_{\|x\|=1} (Ax, x)$
\begin{theorem}[11.5]\label{11.5}
    $H(\C)$, $A$~--- самосопряжённый оператор.
   \begin{enumerate}
        \item $\sigma(A) \subset [m_-, m_+]$, 
        причём $m_-, m_+ \in \sigma(A)$, то есть
        $r(A) = \max(|m_-|, |m_+|)$
        \item $r(A)=\|A\|$
    \end{enumerate}
\end{theorem}


\begin{proof}
\begin{itemize}
      \item[2.] $r(A)=\lim_{n \to \infty} \sqrt[n]{\|A^n\|}$. Тогда остаётся доказать лемму:
        \begin{lemma}
        Для самосопряжённого оператора $\|A^n\| = \|A\|^n$
        \end{lemma}
        % мы умеем говорить
        \textbf{Доказательство леммы:} Достаточно понять, почему $\|A^2\| = \|A\|^2$. Мы знаем, что $\|A^2\| \leqslant \|A\|^2$.
        \[
            \|A^2 x\| \cdot \|x\| \geqslant (A^2 x, x) = (Ax, Ax)
        \]
        Возьмем супремум от обеих частей
        \[
            \|A^2\| = \sup_{\|x\|=1} \|A^2 x\| \cdot \| x \| = \sup_{\|x\|=1} \|Ax\|^2=\|A\|^2
        \]
        Получили, что $\|A^2\| \geqslant \|A\|^2$.
        % прыгаем как кенгуру
        В общем-то, пока мы доказали лемму только для степеней двойки (и этого достаточно для доказательства теоремы).
        Абсолютно аналогично можно показать, что лемма верна для всех $n$.
        Переход от $n$ к $n+1$ остается в качестве упражнения.

        \item[1.] 
        \begin{itemize}
            \item[a)] $\sigma(A) \subset [m_-, m_+]$
            по теореме \nameref{11.4} спектр лежит на вещественной оси. Из пункта 2 он лежит на отрезке $[-\|A\|; \|A\|]$. 
            Достаточно показать, что если возьмем число больше $m_+$ или меньше $m_-$, то попадем в резольвентное множество.
            % тайные знания про параболы

            Покажем, что если $\lambda > m_+ \implies \lambda \in \rho(A)$. Будем доказывать по теореме \nameref{11.3}:
            нужно получить оценку вида $\|A_{\lambda} x\| \geqslant m \|x\|$.

            $m_+ = \sup_{\|x\| = 1} (Ax, x)$

            \[
            \|A_\lambda x \| \|x\| \geqslant |(A_{\lambda} x, x)| = |\underbrace{(A x, x)}_{\leqslant m_+ \|x\|^2} - \underbrace{\lambda}_{> m_+} (x, x)| = 
            \lambda (x, x) - (A x, x) \geqslant (\lambda - m_+) \|x\|^2
            \]

            Сокращая на $\|x\|$, получаем:
            $\|A_{\lambda} x\| \cdot \|x\| \geqslant (\lambda - m_+) \|x\|$.
            % хотим попасть в цирк
            \item[b)] Покажем, что $m_{+}$ является частью спектра. Хотим воспользоваться теоремой \nameref{11.3} (пункт 2): существует последовательность $x_n$, такая что $\|x_n\| = 1 \; \forall n$ и $\|A_{\lambda} x_n\| \to 0$.

            Возьмем $x_n$: $Ax_n \to m_+, \: \|x_n\| = 1$ (так как $m_+$ это супремум). 

            Можно переписать это так:
            
            $(A x_n, x_n) - m_+ (x_n, x_n) \to 0$

            $(A_{m_+} x_n, x_n) \to 0$

            Обозначим через $B$ оператор $-(A - m_+ I)$ (он ограниченный и самосопряжённый).

            $(Bx, x) \geqslant 0$

            Введём полускалярное произведение ($\angles{x, x}$ может быть равен 0 при ненулевом $x$): $\angles{x, y} = (Bx, y)$.

            \begin{statement}
                Неравенство Коши-Буняковского справедливо для полускалярного произведения.
                
            \end{statement}
            \textbf{Доказательство утверждения (идейно):} Доказательство оставляем читателю в качестве упражнения.

                Подсказки: 

                Как мы доказываем обычное неравенство Коши-Буняковского?

                \[
                |\angles{x, e}| \leqslant \|x\|
                \| x - (x, e) e\|^2 = \|x\|^2 - |(x, e)|^2 \geqslant 0
                \]

                В чём проблема?

                Если $\angles{x, x} \neq 0$ или $\angles{y, y} \neq = 0$, тогда работает старая схема. Остаётся рассмотреть случай, когда $\angles{x, x} = \angles{y, y} = 0$.
                В вещественном случае рассмотрим выражение 
                \[
                \angles{x \pm y, x \pm y} = 0 + 2 \angles{x, y} + 0
                \geqslant 0
                \]

                В комплексном случае для действительной части аналогично, а для мнимой возьмем $x+iy$, чтобы добраться до мнимой части $\blacksquare$
            
            
            Рассмотрим $\|Bx_n\|^4$: % поставим пока точку
            % нормируем пальцы

            \[
            \|Bx_n\|^4 = (Bx_n, Bx_n)^2 = \angles{x_n, B x_n}^2 \stackrel{\text{КБ}}{\leqslant}  
            \underbrace{\angles{x_n, x_n}}_{=(B x_n, x_n) \to 0} 
            \cdot \underbrace{\angles{B x_n, B x_n}}_{(B^2 x_n, B x_n)
            \stackrel{\text{КБ}}{\leqslant} \|B\|^3 \|x_n\|^2 = 
            \|B\|^3} \to 0
            \]

            % стираем мелом
            % Грам точно не из Италии (он из Дании)
            % посылка с неравенством Йенсена
        \end{itemize}
  \end{itemize}
\end{proof}